\subsection{The last four profiles in the seventeen-arrow census}

\begin{proposition}[Proved: three pendant small-corner profiles]
\label{CP2:A-four17-final-pendant}
None of the following four-object hom-cardinality matrices
has a classifying space of type $\CP^2$:
\[
\begin{gathered}
\begin{pmatrix}1&0&0&1\\0&1&0&1\\0&0&7&3\\0&0&0&3\end{pmatrix},\qquad
\begin{pmatrix}1&0&0&1\\0&6&2&3\\0&0&1&0\\0&0&0&3\end{pmatrix},\\
\begin{pmatrix}1&0&0&1\\0&7&2&2\\0&0&1&0\\0&0&0&3\end{pmatrix}.
\end{gathered}
\]
These are rows $5,237,255$ of the frozen four-object
seventeen-arrow residual ledger, respectively.
\end{proposition}
\begin{proof}
Let $N$ be the corner at object $3$, of order three.
It is a sink corner. Object $0$ is trivial and is joined
only to $3$, by a singleton stalk. Let $D$ be the full
deletion of object $0$. The full subcategory on $0,3$
has initial object $0$, hence contractible classifying
space. Its nerve and $ND$ cover $NC$, and their intersection
is the nerve of $N$: no strict chain can pass through $0$
and a further object other than $3$.

First suppose $BN$ is contractible. The inclusions of
these nerves are CW cofibrations, so their union is their
homotopy pushout. The actual inclusion $BD\to BC$ is
therefore a homotopy equivalence. Deleting $0$ removes
exactly its identity and its singleton arrow, so $D$ has
fifteen arrows. The complete sixteen-arrow exclusion
rules out $BD\simeq\CP^2$.

Otherwise the small-corner classification
Corollary~\ref{CP2:A-five-types} identifies $BN$ with the
classifying space of a nontrivial finite minimal group.
The minimal ideal of $N$ is that group itself. Indeed,
a nontrivial group with two rows or two columns would
already require at least four ideal elements; if the
ideal is not a group, an identity outside it is additionally
required. This is impossible at order three. Averaging
over the group ideal makes both trivial rational $N$
modules projective.

Apply the actual sink-orbit functor
$C\to\overline C$ from Lemma~\ref{CP2:A-high-orbits}:
replace $N$ by the trivial corner and every incoming
stalk by its $N$-orbit set. It induces a rational
cohomology-ring isomorphism. All the stalks remain
nonempty, so the comparability support is unchanged.
In each of the three displayed profiles that support is
a tree. After the orbit functor there is just one corner
of order greater than two, and its order is six or seven.
Under the assumption $BC\simeq\CP^2$, the space
$B\overline C$ has bounded rational homology.
Proposition~\ref{CP2:A-tree-one-large-corner-eleven}
annihilates all its degree-two products. Pullback through
the specified functor contradicts the nonzero rational
square on $BC$. This proves all three exclusions.
\end{proof}

\begin{proposition}[Proved: the final triangle with a small sink]
\label{CP2:A-four17-final-triangle}
The hom-cardinality matrix
\[
\begin{pmatrix}
1&0&0&1\\
0&6&2&3\\
0&0&1&1\\
0&0&0&2
\end{pmatrix}
\]
cannot have classifying space of type $\CP^2$.
It is row $238$ of the frozen residual ledger.
\end{proposition}
\begin{proof}
The sink corner at object $3$ has order two, so its
trivial rational modules are projective on both sides.
The actual sink-orbit functor $C\to\overline C$ replaces
that corner by the trivial monoid and gives a rational
cohomology-ring isomorphism. Its singleton stalks from
$0$ and $2$ remain singletons. The stalk $1\to2$ still
has two arrows, and the stalk $1\to3$ becomes a nonempty
act of the six-element corner at $1$.

Delete $0$ from $\overline C$. Its only outgoing cross
arrow is the singleton arrow to the trivial corner $3$.
Lemma~\ref{CP2:A-five18-singleton-deletion}, in its dual
form, proves that the full deletion inclusion is a
homotopy equivalence. In the resulting category on
$1,2,3$, delete $2$ by the same argument: its only
outgoing cross arrow is again the singleton to $3$.
The resulting full category $E$ on $1,3$ is an ordinary
act cone over the six-element corner, and the two actual
inclusions give $BE\simeq B\overline C$.

If $BC$ were a projective plane, this comparison would
give bounded rational homology of $BE$.
Theorem~\ref{CP2:A-bounded-eleven-acts} annihilates its
degree-two products. The two homotopy equivalences and
the orbit functor transfer this vanishing to $BC$, a
contradiction. The argument uses rational boundedness
after the orbit functor, not an asserted integral
homotopy equivalence for that functor.
\end{proof}
